%%%%%%%%%%%%%%%%%%%%%%%%%%%%%%%%%%%%%%%%%%%%%%%%%%%%%%%%%%%%%%%%%%%%%%%%%%%
% Plantilla genérica para TFG/TFM/Tesis.
%
% Función: configuración del documento y comandos personalizados.
%
% Uso: debes modificar este fichero para adaptarlo a tu trabajo.
% Define la titulación, el título, tus datos, los tutores, el idioma,
% la estructura y las opciones que quieras utilizar.
% Consulta los valores admitidos en los comentarios de cada variable.
% Al final puedes añadir tus propios comandos.
%
% Estos datos se reutilizan en el documento, las páginas institucionales
% y los formularios, para evitar que tengas que introducirlos varias veces.
%%%%%%%%%%%%%%%%%%%%%%%%%%%%%%%%%%%%%%%%%%%%%%%%%%%%%%%%%%%%%%%%%%%%%%%%%%%

%%%%%%%%%%%%%%%%%%%%%%%%%%%%%%%%%%%%%%%%%%%%%%%%%%%%%%%%%%%%%%%%%%%%%%%%%%% 
% INICIO Configuración de las variantes del documento

% Selecciona el tipo de documento (TFM, TFG, tesis, etc.).
% La titulación puede tomar los siguientes valores:
% Los identificadores ajenos a la UAH usan <degree>-<university>; conserva el guion.
% Grados de la UAH:
%  En la Escuela Politécnica Superior:
%   GITT   - Grado en Ingeniería en Tecnologías de la Telecomunicación
%   GIEC   - Grado en Ingeniería Electrónica de Comunicaciones
%   GIT    - Grado en Ingeniería Telemática
%   GIST   - Grado en Ingeniería en Sistemas de Telecomunicación
%   GIC    - Grado en Ingeniería de Computadores
%   GII    - Grado en Ingeniería Informática
%   GMC    - Grado en Matemáticas y Computación
%   GSI    - Grado en Sistemas de Información
%   GISI   - Grado en Ingeniería en Sistemas de Información
%   GIEAI  - Grado en Ingeniería en Electrónica y Automática Industrial
%   GITI   - Grado en Ingeniería en Tecnologías Industriales
%  En la Facultad de Ciencias:
%   GFIE   - Grado en Física e Instrumentación Espacial
%   GBS    - Grado en Biología Sanitaria
%   GB     - Grado en Biología
%   GCA    - Grado en Ciencias Ambientales
%   GQ     - Grado en Química
%   GCCTF  - Grado en Criminalística: Ciencias y Tecnologías Forenses
%
% Másteres de la UAH:
%  En la Escuela Politécnica Superior:
%   MUSEA  - Máster Universitario en Sistemas Electrónicos Avanzados. Sistemas Inteligentes
%   MUIT   - Máster Universitario en Ingeniería de Telecomunicación
%   MUII   - Máster Universitario en Ingeniería Industrial
%   MUIE   - Máster Universitario en Ingeniería Electrónica
%   MUCTE  - Máster Universitario en Ciencia y Tecnología desde el Espacio
%   MUC    - Máster Universitario en Ciberseguridad
%   MUANBD - Máster Universitario en Analítica de Negocio y Big Data
%
% Programas de doctorado:
%  PHDUAH  - Doctorado UAH
%  PHD-ETSIT-UPM - Doctorado de la ETSIT-UPM
%
% Titulaciones de la URJC:
%  ITI-URJC - Grado en Ingeniería en Tecnologías Industriales (ESCET-URJC)
%
% Titulaciones de la ETSIT-UPM:
%  GISTS-UPM  - Grado en Ingeniería de Tecnologías y Servicios de Telecomunicación
%  GIB-UPM    - Grado en Ingeniería Biomédica
%  GISD-UPM   - Grado en Ingeniería y Sistemas de Datos
%  MUIT-UPM   - Máster Universitario en Ingeniería de Telecomunicación
%  MUCS-UPM   - Máster Universitario en Ciberseguridad
%  MUESFV-UPM - Máster Universitario en Energía Solar Fotovoltaica
%  MUIB-UPM   - Máster Universitario en Ingeniería Biomédica
%  MUIRST-UPM - Máster Universitario en Ingeniería de Redes y Servicios Telemáticos
%  MUISE-UPM  - Máster Universitario en Ingeniería de Sistemas Electrónicos
%  MUTSC-UPM  - Máster Universitario en Teoría de la Señal y Comunicaciones
%  MUNE-UPM   - Máster Universitario en Neurotecnología (páginas en español o inglés)
%
% Otros documentos:
%  GEINTRARR - Geintra Research Report (soporte preliminar)
%
% Especifica aquí tu titulación, de entre las listadas arriba.
\newcommand{\myDegree}{GISI}

\newcommand{\mySpecialty}{} % Grados de la UAH y perfil/especialidad en la hoja de control del anteproyecto de TFM EPS-UAH.

% Selecciona el idioma del documento y sus ajustes específicos.
% Los valores admitidos son english o spanish.
\newcommand{\myLanguage}{spanish}

% Selecciona la estructura del documento:
%  standard   - estructura habitual para TFG, TFM, tesis convencionales y
%               los demás documentos compatibles con la plantilla.
%               Organiza su contenido en Book/content-standard.tex.
%  compendium - estructura específica para tesis doctorales presentadas
%               como compendio de publicaciones. Organiza su contenido en
%               Book/content-compendium.tex.
\newcommand{\myDocumentStructure}{standard}

% Estilo visual del documento principal; no afecta al papeleo:
% standard  - aspecto original de la plantilla (opción predeterminada recomendada)
% editorial - diseño con serifas y un separador vertical discreto en los capítulos
% modern    - títulos sin serifas, color para resaltar elementos y números de capítulo grandes
% framed    - títulos de capítulo enmarcados, al estilo Lenny
% shaded    - títulos de capítulo sombreados, al estilo Bjornstrup
% mimosis   - diseño basado en Garamond e inspirado en la plantilla Mimosis
\newcommand{\myTypesettingStyle}{standard}



%%%%%%%%%%%%%%%%%%%%%%%%%%%%%%%%%%%%%%%%%%%%%%%%%%%%%%%%%%%%%%%%%%%%%%%%%%%
% Información general del documento
\newcommand{\myBookTitleSpanish}{Sistema de Generación Automática de Setlists para DJs a partir del Análisis Musical de una Biblioteca de Audio}
\newcommand{\myBookTitleEnglish}{Automatic DJ Setlist Generation System Based on Musical Analysis of an Audio Library}
\newcommand{\myThesisKeywords}{Sistema de recomendación, mezcla armónica, Rekordbox, DJ, automatización} % (máximo de cinco)
\newcommand{\myThesisKeywordsEnglish}{Recommendation system, harmonic mixing, Rekordbox, DJ, automation} % (máximo de cinco)
\newcommand{\myConfidentialContent}{No} % Puede ser Yes o No; se utiliza
% en el papeleo del MUII desde julio de 2021.

%%%%%%%%%%%%%%%%%%%%%%%%%%%%%%%%%%%%%%%%%%%%%%%%%%%%%%%%%%%%%%%%%%%%%%%%%%%
% Datos del autor
\newcommand{\myAuthorName}{Pablo}
\newcommand{\myAuthorSurname}{Hernández Borregón}
\newcommand{\myAuthorFullName}{\myAuthorName{} \myAuthorSurname{}}
\newcommand{\myAuthorGender}{male} 
\newcommand{\myAuthorEmail}{pablo.hernandezb@edu.uah.es}
\newcommand{\myAuthorDNI}{09823707Q}
% Datos personales para la solicitud de anteproyecto; no siempre son necesarios.
\newcommand{\myAuthorStreet}{C/: Senda de la Cultura}
\newcommand{\myAuthorCity}{Alcalá de Henares}
\newcommand{\myAuthorPostalCode}{28806}
\newcommand{\myAuthorProvince}{Madrid}
\newcommand{\myAuthorTelephone}{637689756}


%%%%%%%%%%%%%%%%%%%%%%%%%%%%%%%%%%%%%%%%%%%%%%%%%%%%%%%%%%%%%%%%%%%%%%%%%%%
% Datos del tutor
% GISD-UPM: el tutor definido en myAcademicTutorFullName se muestra como
% supervisor técnico, no como supervisor académico.
\newcommand{\myAcademicTutorFullName}{Óscar Gutiérrez Blanco} 
\newcommand{\myAcademicTutorGender}{male}
\newcommand{\myAcademicTutorDNI}{}
% Departamento e institución son campos separados; deja vacío lo que no proceda.
\newcommand{\myAcademicTutorDepartment}{Ciencias de la Computación}
\newcommand{\myAcademicTutorDepartmentEnglish}{Computer Science}
\newcommand{\myAcademicTutorInstitution}{Universidad de Alcalá}
\newcommand{\myAcademicTutorInstitutionEnglish}{University of Alcalá}
\newcommand{\myAcademicTutorPosition}{T.U. de CC. de Computación e Int. Art.}
\newcommand{\myAcademicTutorPositionEnglish}{Professor of Computer Science and Artificial Intelligence} % Sólo para titulaciones UPM: créditos de tesis en inglés.
\newcommand{\myAcademicTutorEmail}{oscar.gutierrez@uah.es}
\newcommand{\myAcademicTutorForeignUniversity}{}
\newcommand{\myAcademicTutorForeignCity}{}
\newcommand{\myAcademicTutorIsPhD}{true} % true o false, según la normativa de TFM de la EPS-UAH.

% Sólo para titulaciones UPM: Ponente ETSIT-UPM (GISTS/GIB y másteres con su
% mismo modelo): papel distinto del tutor y del cotutor. Déjalo vacío si todavía
% no está asignado; la hoja del tribunal deja el valor vacío.
\newcommand{\myPonenteFullName}{}
% Modelo GISTS/GIB y másteres asociados: si hay ponente, se muestra su
% departamento; si no, el del tutor.
\newcommand{\myPonenteDepartment}{}

% Sólo para titulaciones UPM: Supervisión técnica del MUNE-UPM, distinta de la
% supervisión académica (myAcademicTutorFullName). Puede ser la misma persona;
% indica ambos nombres. Si todavía no está asignada, déjalo vacío para mostrar
% una línea en blanco.
\newcommand{\myUPMTechnicalSupervisorFullName}{}
% GISD-UPM no utiliza myUPMTechnicalSupervisorFullName: usa el tutor y cotutor.

%%%%%%%%%%%%%%%%%%%%%%%%%%%%%%%%%%%%%%%%%%%%%%%%%%%%%%%%%%%%%%%%%%%%%%%%%%%
% Datos del cotutor
% GISD-UPM: el cotutor definido en myCoTutorFullName se muestra como supervisor
% académico; completa este campo y myAcademicCotutorDepartment.
\newcommand{\myCoTutorFullName}{Ángel Javier Álvarez Miguel} % Déjalo vacío si no hay cotutor.
\newcommand{\myCoTutorGender}{male}
\newcommand{\myCoTutorDNI}{}
\newcommand{\myAcademicCotutorDepartment}{Servicio de Aplicaciones}
\newcommand{\myAcademicCotutorDepartmentEnglish}{Application Services}
\newcommand{\myAcademicCotutorInstitution}{Universidad de Alcalá}
\newcommand{\myAcademicCotutorInstitutionEnglish}{University of Alcalá}
\newcommand{\myCoTutorPosition}{Técnico/a Informático/a - Gestión Administrativa}
\newcommand{\myCoTutorPositionEnglish}{IT Technician - Administrative Management} % Sólo para titulaciones UPM: créditos de tesis en inglés.
\newcommand{\myCoTutorEmail}{a.alvarez@uah.es}
\newcommand{\myCoTutorIsPhD}{true} % true o false, según la normativa de TFM de la EPS-UAH.

%%%%%%%%%%%%%%%%%%%%%%%%%%%%%%%%%%%%%%%%%%%%%%%%%%%%%%%%%%%%%%%%%%%%%%%%%%%
% Datos institucionales
% La universidad, su acrónimo y el centro se seleccionan automáticamente en
% Config/degrees.tex y Config/institutions.tex según \myDegree.

% El departamento del documento se obtiene del tutor o, si está vacío, del
% cotutor. Configúralo en sus campos Department/DepartmentEnglish anteriores.

\newcommand{\myPhDProgram}{}
\newcommand{\myPhDProgramEnglish}{}

% Sólo para titulaciones UPM: Tesis PHD-ETSIT-UPM: utiliza el nombre oficial
% completo del programa en los campos anteriores (Doctorado en ... / Doctoral
% Degree in ...), sin abreviarlo. Titulación previa de acceso al doctorado, no
% el título de doctor que solicitas.
\newcommand{\myUPMdegree}{}
\newcommand{\myUPMdegreeEnglish}{}
% Año de solicitud del depósito; puede ser distinto del año de defensa.
\newcommand{\myThesisDepositYear}{\myThesisDefenseYear}
% Texto completo de financiación competitiva para los créditos PHD-ETSIT-UPM.
% Déjalo vacío si no procede; admite texto LaTeX y no aparece en otras titulaciones.
\newcommand{\myThesisFundingStatement}{}
\newcommand{\myThesisFundingStatementEnglish}{}

\newcommand{\myResearchGroup}{}


%%%%%%%%%%%%%%%%%%%%%%%%%%%%%%%%%%%%%%%%%%%%%%%%%%%%%%%%%%%%%%%%%%%%%%%%%%%
% Miembros del tribunal y personal del departamento
\newcommand{\myTribunalPresident}{}
\newcommand{\myTribunalPresidentGender}{male}
\newcommand{\myTribunalPresidentDNI}{} % Necesario para el compromiso de confidencialidad.
% ETSIT-UPM: el primer vocal (FirstSpokesperson) se muestra como Vocal.
\newcommand{\myTribunalFirstSpokesperson}{}
\newcommand{\myTribunalFirstSpokespersonGender}{male}
\newcommand{\myTribunalFirstSpokespersonDNI}{} % Necesario para el compromiso de confidencialidad.
\newcommand{\myTribunalSecondSpokesperson}{}
\newcommand{\myTribunalSecondSpokespersonGender}{female}
\newcommand{\myTribunalSecondSpokespersonDNI}{} % Necesario para el compromiso de confidencialidad.
\newcommand{\myTribunalAlternateMember}{}
\newcommand{\myTribunalAlternateMemberGender}{male}
\newcommand{\myTribunalAlternateMemberDNI}{} % Necesario para el compromiso de confidencialidad.
\newcommand{\myTribunalSecretary}{}
\newcommand{\myTribunalSecretaryGender}{male} % male o female; independiente de la secretaría del departamento
\newcommand{\myDepartmentSecretary}{} % Para el papeleo de TFG y TFM, incluido MUSEA.
\newcommand{\myDepartmentSecretaryGender}{female}                  % Para el papeleo de TFG y TFM, incluido MUSEA.
\newcommand{\myTFMComisionPresident}{} % Nombre utilizado en el papeleo del MUIE.
\newcommand{\myTFMComisionPresidentGender}{male} % male o female; cargo en la hoja de control del anteproyecto de TFM EPS-UAH, independiente del tribunal.

%%%%%%%%%%%%%%%%%%%%%%%%%%%%%%%%%%%%%%%%%%%%%%%%%%%%%%%%%%%%%%%%%%%%%%%%%%%
% Sólo para titulaciones UPM: Tribunal ampliado y revisores externos para tesis
% PHD-ETSIT-UPM. Completa los nombres asignados; deja vacíos los campos que no
% procedan.
\newcommand{\myTribunalThirdSpokesperson}{}
\newcommand{\myTribunalSecondAlternateMember}{}
\newcommand{\myFirstExternalReviewerFullName}{}
\newcommand{\myFirstExternalReviewerInstitution}{}
\newcommand{\myFirstExternalReviewerInstitutionEnglish}{}
\newcommand{\mySecondExternalReviewerFullName}{}
\newcommand{\mySecondExternalReviewerInstitution}{}
\newcommand{\mySecondExternalReviewerInstitutionEnglish}{}

%%%%%%%%%%%%%%%%%%%%%%%%%%%%%%%%%%%%%%%%%%%%%%%%%%%%%%%%%%%%%%%%%%%%%%%%%%%
% Calificaciones. Si las defines, la nota final se calcula automágicamente
% en el documento de acta y rúbrica.
\newcommand{\myAcademicTutorGrade}{}
\newcommand{\myTribunalGrade}{}
\newcommand{\myTribunalMHProposal}{NO} % Los valores admitidos son YES o NO.


%%%%%%%%%%%%%%%%%%%%%%%%%%%%%%%%%%%%%%%%%%%%%%%%%%%%%%%%%%%%%%%%%%%%%%%%%%%
% Fechas
\newcommand{\myThesisProposalDate}{} % Fecha del anteproyecto.

\newcommand{\myThesisDepositDate}{}
\newcommand{\myThesisDepositDateEnglish}{}

% En los RR, myThesisDefenseDate es la fecha que aparece en la portada.
\newcommand{\myThesisAcademicYear}{2026/2027}
\newcommand{\myThesisDefenseYear}{2027}
\newcommand{\myThesisDefenseDate}{}
\newcommand{\myThesisdefenseDateEnglish}{}
% Si prefieres el formato británico para la fecha en inglés, usa este ejemplo:
% \newcommand{\myThesisdefenseDateEnglish}{6\textsuperscript{th} of January, 2018}

\newcommand{\myPaperworkDate}{\myThesisDepositDate}
\newcommand{\myPaperworkDateEnglish}{\myThesisDepositDateEnglish}

%%%%%%%%%%%%%%%%%%%%%%%%%%%%%%%%%%%%%%%%%%%%%%%%%%%%%%%%%%%%%%%%%%%%%%%%%%%
% Datos para la publicación en abierto
\newcommand{\myAuthorizationOpenPublishing}{Yes}
\newcommand{\myAuthorizationOpenPublishingEmbargoMonths}{0} % Puede ser
%  0,  6, 12,
% 18, 24

\newcommand{\myResearchVicerrector}{Excmo. Sr. Francisco J. de la Mata de la Mata}
\newcommand{\myResearchVicerrectorGender}{male}

% Derechos de autor
\newcommand{\myCopyrightStatement}{\myAuthorFullName{}. Some rights reserved. This document is under terms of Creative Commons license Attribution - Non Commercial - Non Derivatives.}
\newcommand{\myLicenseURL}{http://creativecommons.org/licenses/by-nc-nd/3.0/es/}

\newcommand{\myResearchReportID}{RR-2027-01}

%%%%%%%%%%%%%%%%%%%%%%%%%%%%%%%%%%%%%%%%%%%%%%%%%%%%%%%%%%%%%%%%%%%%%%%%%%%
% Elementos opcionales del documento. Cada valor debe ser exactamente true o false.
% Desactivar una lista no desactiva la función LaTeX correspondiente.
\newcommand{\myIncludeSampleLetter}{false}
\newcommand{\myIncludeDedication}{false}
\newcommand{\myIncludeAcknowledgements}{true}
\newcommand{\myIncludeListOfFigures}{true}
\newcommand{\myIncludeListOfTables}{true}
\newcommand{\myIncludeListOfAcronyms}{true}
\newcommand{\myIncludeListOfSymbols}{false}
\newcommand{\myIncludeListOfCode}{false}
\newcommand{\myIncludeListOfAlgorithms}{false}
\newcommand{\myIncludeListOfVideos}{false}

\newcommand{\myFlagSplittedAdvisors}{true} % Si es false, muestra
% "Tutores/Advisors" en las portadas.
% Si es true, separa Tutor/Cotutor
% o Advisor/Co-advisor. En MUIT y MUII
% se ignora y se fija en true.

%%%%%%%%%%%%%%%%%%%%%%%%%%%%%%%%%%%%%%%%%%%%%%%%%%%%%%%%%%%%%%%%%%%%%%%%%%%
% Colores de los enlaces
% Enlaces de los índices general, de tablas y de figuras.
%\newcommand{\mytoclinkcolor}{blue}
\newcommand{\mytoclinkcolor}{black}
%\newcommand{\myloflinkcolor}{red}
\newcommand{\myloflinkcolor}{black}
%\newcommand{\mylotlinkcolor}{green}
\newcommand{\mylotlinkcolor}{black}

% Enlaces de las listas de código, algoritmos y vídeos.
%\newcommand{\myothertoclinkcolor}{magenta}
\newcommand{\myothertoclinkcolor}{black}

% Otros enlaces internos del documento.
\newcommand{\mylinkcolor}{blue}
%\newcommand{\mylinkcolor}{black}

% Enlaces a URL y citas bibliográficas.
\newcommand{\myurlcolor}{blue}
%\newcommand{\myurlcolor}{black}
%\newcommand{\mycitecolor}{green}
\newcommand{\mycitecolor}{blue}
%\newcommand{\mycitecolor}{black}

% FIN Configuración de las variantes del documento
%%%%%%%%%%%%%%%%%%%%%%%%%%%%%%%%%%%%%%%%%%%%%%%%%%%%%%%%%%%%%%%%%%%%%%%%%%% 

%%%%%%%%%%%%%%%%%%%%%%%%%%%%%%%%%%%%%%%%%%%%%%%%%%%%%%%%%%%%%%%%%%%%%%%%%%% 
% INICIO Comandos personalizados
% Define aquí tus propios comandos.

% Este comando define el formato del texto en inglés dentro de un documento
% en español.
\DeclareRobustCommand{\texten}[1]{\textit{#1}}

\def\ci{\perp\!\!\!\perp}

% Ejemplos de comandos de uso frecuente.
\newcommand{\circulo}{\large $\circ$}
\newcommand{\asterisco}{$\ast$}
\newcommand{\cuadrado}{\tiny $\square$}
\newcommand{\triangulo}{\scriptsize $\vartriangle$}
\newcommand{\triangv}{\scriptsize $\triangledown$}
\newcommand{\diamante}{\large $\diamond$}

\newcommand{\new}[1]{\textcolor{magenta}{#1 }}
\newcommand{\argmax}[1]{\underset{#1}{\operatorname{argmax}}}

% Estos comandos se utilizan en los capítulos de ejemplo.
\newcommand{\verticalSpacingSRPMaps}{-0.3cm}

\newcommand{\CAVDB}{\texttt{CAV3D}}
\newcommand{\AVThreeTVLoc}{\texttt{AV3T-VLoc}}
\newcommand{\LandPoseVLoc}{\texttt{LandPose-VLoc}}
\newcommand{\DeltaPoseAVT}{$\Delta_{\AVThreeTVLoc}^{\LandPoseVLoc}$}
\newcommand{\mae}{\epsilon{}}
\newcommand{\maeThreeD}{\mae_{3D}}
\newcommand{\maeThreeDfine}{\maeThreeD'}
\newcommand{\maeTwoD}{\mae_{2D}}
\newcommand{\maeTwoDfine}{\maeTwoD'}

\newcommand{\NumDataFrames}{N_F}


% Anotaciones de revisión
\newcommand{\todojmg}{\todo[author=JMG, inline, color=green!20, caption={}]}
\newcommand{\todorbc}{\todo[author=RBC, inline, color=red!20,   caption={}]}
\newcommand{\todomig}{\todo[author=MIG, inline, color=blue!20,  caption={}]}
% \newcommand{\todospc}{\todo[author=SPC, inline, color=green!60,  caption={}]}
% \newcommand{\todolbm}{\todo[author=LBM, inline, color=green!40, caption={}]}
% \newcommand{\todojmb}{\todo[author=JMB, inline, color=red!40,   caption={}]}
% \newcommand{\todoapt}{\todo[author=APT, inline, color=blue!40,  caption={}]}

% FIN Comandos personalizados
%%%%%%%%%%%%%%%%%%%%%%%%%%%%%%%%%%%%%%%%%%%%%%%%%%%%%%%%%%%%%%%%%%%%%%%%%%% 


%%%%%%%%%%%%%%%%%%%%%%%%%%%%%%%%%%%%%%%%%%%%%%%%%%%%%%%%%%%%%%%%%%%%%%%%%%%
% INICIO Definiciones obsoletas o de menor utilidad; consérvalas por compatibilidad.
\newcommand{\myFirstAdvisorFullName}{\myAcademicTutorFullName} % Obsoleta: toma el nombre del tutor.
\newcommand{\mySecondAdvisorFullName}{\myCoTutorFullName} % Obsoleta: toma el nombre del cotutor.
\newcommand{\myFirstAdvisorDNI}{\myAcademicTutorDNI} % Obsoleta: toma el DNI del tutor.
\newcommand{\mySecondAdvisorDNI}{\myCoTutorDNI} % Obsoleta: toma el DNI del cotutor.
\newcommand{\mybookFigure}{alumno} % Obsoleta: indicaba en los TFG
% la condición del autor que firmaba
% el convenio (normalmente "alumno").

% FIN Definiciones obsoletas o de menor utilidad
%%%%%%%%%%%%%%%%%%%%%%%%%%%%%%%%%%%%%%%%%%%%%%%%%%%%%%%%%%%%%%%%%%%%%%%%%%%

%%% Local Variables:
%%% TeX-master: "../book"
%%% End:
